\editionPage[\pdfbookmark[0]{Comprendre pourquoi}{bm-preuves-debut-parent}]{preuves-debut}{Comprendre pourquoi}{Pourquoi 1 dm² vaut 100 cm²}{Les explications ci-dessous éclairent des règles utilisées dans le livre.}
\begin{blocEdition}{preuves-debut}
\editionEtape{Paver le grand carré}Un carré de 1 dm de côté mesure 10 cm de côté. On peut y ranger une rangée de 10 petits carrés de 1 cm de côté.\editionEtape{Compter les rangées}Il y a 10 rangées identiques, l’une au-dessus de l’autre : $10\times10=100$ petits carrés de $1\ \mathrm{cm}^2$.\editionEtape{Généraliser}De même, 1 m$^2$ contient 10 rangées de 10 dm$^2$, soit $100\ \mathrm{dm}^2$. Pour les aires, le rapport entre deux unités voisines est $10\times10=100$, et non 10.\editionRetenir{Une aire croise deux longueurs. Quand l’unité de longueur est multipliée par 10, l’unité d’aire est multipliée par 100.}
\end{blocEdition}
\editionPage{preuves-2}{Comprendre pourquoi}{Pourquoi on additionne seulement les numérateurs}{Les explications ci-dessous éclairent des règles utilisées dans le livre.}
\begin{blocEdition}{preuves-2}
\editionEtape{Des parts identiques}Sur une bande partagée en huit, chaque part est un huitième : toutes les parts ont la même taille.\label{fig-preuves-2-preuve-fractions}\figureGeo{preuve-fractions}
\editionEtape{Réunir}Trois huitièmes et deux huitièmes font cinq huitièmes, comme 3 billes et 2 billes font 5 billes : $\frac38+\frac28=\frac58$.\editionEtape{Des parts différentes}Pour $\frac34+\frac18$, on rend d’abord les parts identiques : un quart vaut deux huitièmes, donc $\frac34=\frac68$ et la somme vaut $\frac78$.\editionRetenir{Le dénominateur décrit le partage ; il ne s’additionne jamais. On n’ajoute que des parts de même taille.}
\end{blocEdition}
\editionPage{preuves-3}{Comprendre pourquoi}{Pourquoi le reste est plus petit que le diviseur}{Les explications ci-dessous éclairent des règles utilisées dans le livre.}
\begin{blocEdition}{preuves-3}
\editionEtape{Faire des paquets}Diviser 87 par 4, c’est former le plus possible de paquets de 4. Tant qu’il reste au moins 4 objets, on peut encore former un paquet.\editionEtape{Quand s’arrêter}On s’arrête quand il reste 0, 1, 2 ou 3 objets : $87=4\times21+3$. Un reste de 4 ou plus voudrait dire qu’un paquet a été oublié.\editionEtape{Contrôler une division}Un résultat se vérifie donc de deux façons : quotient fois diviseur plus reste redonne le dividende, et le reste est plus petit que le diviseur.\editionRetenir{Si le reste dépasse le diviseur, la division n’est pas terminée : il faut augmenter le quotient.}
\end{blocEdition}
\editionPage{preuves-4}{Comprendre pourquoi}{Pourquoi l’hexagone se ferme exactement}{Les explications ci-dessous éclairent des règles utilisées dans le livre.}
\begin{blocEdition}{preuves-4}
\editionEtape{Des triangles équilatéraux}En reportant le rayon sur le cercle, chaque côté de l’hexagone a la longueur du rayon. Avec le centre, chaque côté forme un triangle dont les trois côtés sont égaux : un triangle équilatéral.\editionEtape{Faire le tour du centre}Six triangles équilatéraux identiques se rangent côte à côte autour du centre et remplissent exactement le tour complet : leurs six angles au centre sont égaux.\editionEtape{Conclure}Après six reports du rayon, on revient au point de départ : l’hexagone se ferme. Ses six côtés sont égaux et son périmètre vaut six rayons.\editionRetenir{L’hexagone régulier est formé de six triangles équilatéraux : c’est pour cela que les abeilles pavent leurs rayons sans laisser de trou.}
\end{blocEdition}
